% Sources verifiees WebSearch 2026-07-15 : tension de court-circuit Ucc typique 4-6% pour
% transformateurs de distribution MT/BT <= 20 kV (electrical-installation.org, se.com FAQ Schneider)
% Correction 2026-08-25 (audit AUDIT_EXERCICES_2026.md) : exercice "essai a vide" remis sur les
% valeurs authentiques du sujet E2 BAC+3 2026 (V = 230 V, I0 = 42,5 A, P0 = 7,82 kW, questions
% A12-A13, verifiees dans le PDF original) ; exercices reordonnes par difficulte croissante.
% Passe "complementarite fiche <-> exercices" du 2026-09-02 : les deux exercices qui
% reproduisaient les Questions Types de la fiche T06 ont ete remplaces (voir les commentaires
% detailles en tete des exercices 2 et 4).
%   - ancien Ex. 2 (2026 E1, partie E, questions E4 et E10 : P3 vs PD3, m = 2,23 / 1,12)
%     = doublon exact des QT1 et QT2 de la fiche, et empietement sur la notion T08
%     -> remplace par le SEUL exercice docTech de la notion : 2018, EDE ENE 2, questions 13
%        a 15, choix du calibre d'un regulateur abaisseur COMPACTO II dans la gamme
%        constructeur (documents ressources DRS 4 section 8 et DRS 6).
%   - ancien Ex. 4 (2026 E2 BAC+3, questions A12-A13 : essai a vide, Rf = 20,3 ohms,
%     Xm = 5,62 ohms) = doublon exact de la QT3 de la fiche (meme source, memes valeurs)
%     -> remplace par un exercice inedit de dimensionnement du circuit magnetique
%        (formule de Boucherot), travail intellectuel absent des trois QT.
% Les exercices 1, 3 et 5 (inedits : chute de tension par epsilon_cc, couplages Yy/Yd,
% rendement et taux de charge optimal) sont disjoints des QT et n'ont pas ete touches.
\expandafter\ifx\csname CAPETINCLUDEMODE\endcsname\relax
\documentclass[11pt,a4paper]{article}
% preamble.tex — packages et commandes communes a toutes les fiches CAPET SII IE
\input{glyphtounicode}
\pdfgentounicode=1
\usepackage{cmap}
\usepackage[T1]{fontenc}
\usepackage[utf8]{inputenc}
\usepackage[french]{babel}
% babel[french] active a \begin{document} des raccourcis de ponctuation (:;!?)
% qui cassent la syntaxe to[...] de circuitikz (erreur "+ or - expected") ;
% on les desactive, ce qui ne change que l'espacement fin avant ces signes.
\AtBeginDocument{\shorthandoff{:;!?}}
\usepackage[a4paper, top=1.6cm, bottom=1.6cm, left=1.4cm, right=1.4cm]{geometry}
\setlength{\headheight}{14.5pt}
\usepackage{amsmath, amssymb, amsfonts}
\usepackage{siunitx}
\sisetup{locale=FR, output-decimal-marker={,}, inter-unit-product=\ensuremath{{}\cdot{}},
         group-minimum-digits=4}
% Unites non fournies par siunitx, declarees ici une fois pour toutes (fiches + exercices).
% Les gardes \ifdefined evitent tout conflit avec une redeclaration locale restee dans un
% fichier : ne pas retirer les gardes, ne pas redeclarer ces unites dans les fichiers.
\DeclareSIUnit{\jour}{j}
\ifdefined\var\else\DeclareSIUnit{\var}{var}\fi   % puissance reactive
\ifdefined\vad\else\DeclareSIUnit{\vad}{VAD}\fi   % puissance deformante
\ifdefined\tr\else\DeclareSIUnit{\tr}{tr}\fi      % tour (tr/min)
\usepackage{array, booktabs, multirow, tabularx}
\usepackage{xcolor}
\usepackage[most, breakable]{tcolorbox}
\usepackage{multicol}
\usepackage{enumitem}
\usepackage{mathtools}
\usepackage{lmodern}
% Refonte visuelle 2026 (direction "moderne coloré") : Lato pour le texte,
% maths sans-serif assorties via newtxsf (charge apres amsmath/amssymb).
\usepackage[default]{lato}
\usepackage{newtxsf}
\usepackage{microtype}
\usepackage{fancyhdr}
\usepackage{tikz}
\usetikzlibrary{positioning,arrows.meta}
\usepackage{circuitikz}
\usepackage{lastpage}

% ── Couleurs par domaine (palette modernisée 2026) ────────────────────
\definecolor{colElec}{RGB}{29, 78, 216}        % Électrotechnique — bleu vif
\definecolor{colElecLight}{RGB}{219, 234, 254} % fond bleu clair
\definecolor{colAuto}{RGB}{21, 128, 61}        % Automatique — vert
\definecolor{colMeca}{RGB}{194, 65, 12}        % Mécanique — orange
\definecolor{colTherm}{RGB}{185, 28, 28}       % Thermique — rouge
\definecolor{colQT}{RGB}{232, 244, 233}        % (hérité, plus utilisé par questiontype)
\definecolor{colJury}{RGB}{255, 248, 225}      % fond jury (crème)
\definecolor{colJuryFrame}{RGB}{202, 138, 4}   % bordure/bandeau jury (ambre)
\definecolor{colGrayLight}{RGB}{245, 245, 247} % fond gris neutre
\definecolor{colSeparator}{RGB}{180, 180, 180} % gris séparateur
\definecolor{colExo}{RGB}{8, 145, 178}          % Banque d'exercices — cyan (distinct des fiches)
% ── Encre et gris de la refonte ───────────────────────────────────────
\definecolor{colInk}{RGB}{23, 32, 42}          % texte courant (quasi-noir)
\definecolor{colSlate}{RGB}{83, 98, 118}       % gris ardoise (méta, en-têtes)
\definecolor{colRule}{RGB}{203, 213, 225}      % filets clairs
\definecolor{colJuryBg}{RGB}{255, 248, 225}    % fond crème jury
\definecolor{colPresence}{RGB}{83, 98, 118}    % bloc 5 (présence dans les sujets)
\AtBeginDocument{\color{colInk}}

% ── Style encadré exercice + corrigé (banque d'exercices) ─────────────
\tcbset{
  exoBox/.style={
    breakable,
    enhanced,
    colback=colExo!8!white,
    colframe=colExo,
    colbacktitle=colExo,
    coltitle=white,
    fonttitle=\bfseries\small,
    left=4pt, right=4pt, top=3pt, bottom=3pt,
    before skip=5pt, after skip=5pt,
  }
}

% ── Bascule enonce / corrige (banque d'exercices) ─────────────────────
% \CAPETCORRIGE est defini en ligne de commande (pdflatex "\def\CAPETCORRIGE{1}\input{...}")
% pour compiler la variante corriges d'un meme fichier source ; par defaut (non defini),
% la compilation produit la variante enonces. \CAPETSHOW{enonce}{corrige} choisit la branche.
%
% TROIS MODES (refonte 2026-09-20, manuel de revision) :
%   - aucun drapeau   -> imprime l'enonce            (variante enonces, INCHANGE)
%   - \CAPETCORRIGE   -> imprime le corrige          (variante corriges, INCHANGE)
%   - \CAPETMANUEL    -> imprime l'enonce ET met le corrige de cote pour l'annexe
%
% Le mode manuel stocke chaque corrige dans une macro (\csgdef), sans ecriture
% fichier ni re-tokenisation : le contenu (circuitikz, accents, %) traverse intact.
% L'annexe les rejoue ensuite dans l'ordre via \CAPETANNEXE.
\usepackage{etoolbox}

\newcounter{CAPETexo}          % numero continu d'exercice dans le manuel
\setcounter{CAPETexo}{0}
\newcounter{CAPETannexe}       % compteur de relecture en annexe

% Titre courant de la notion en cours (pose par \CAPETNOTION), sert d'intertitre annexe
\newcommand{\CAPETnotioncode}{}
\newcommand{\CAPETnotiontitre}{}

% \CAPETNOTION{code}{titre} : ouvre une notion dans le manuel (memorise le libelle,
% et marque le rang du premier exercice de la notion pour l'intertitre d'annexe).
\newcommand{\CAPETNOTION}[2]{%
  \renewcommand{\CAPETnotioncode}{#1}%
  \renewcommand{\CAPETnotiontitre}{#2}%
  \csgdef{CAPETnotionstart\the\numexpr\value{CAPETexo}+1\relax}{#1 -- #2}%
}

% ── \degree robuste ───────────────────────────────────────────────────
% Certains fichiers d'exercices utilisent \degree. Ni le preambule ni siunitx
% ne le definissent, et \providecommand local ne suffit pas : sous hyperref,
% \degree existe deja comme commande factice (\ERROR), donc \providecommand
% s'abstient et l'appel echoue a la compilation du manuel (erreur invisible
% en compilation isolee). On le (re)definit donc inconditionnellement ici,
% apres le chargement des packages, pour tout le corpus.
\AtBeginDocument{\providecommand{\degree}{}\renewcommand{\degree}{\ensuremath{^\circ}}}

% Libelle decoratif : les fichiers d'exercices utilisent aussi \CAPETSHOW dans
% leur \fancyhead et leur bloc de titre, pour afficher "(enonces)" ou "(corriges)".
% Ces appels-la ne sont PAS des exercices : en mode manuel ils ne doivent ni
% consommer un numero, ni alimenter l'annexe. On les reconnait a leur premier
% argument, strictement egal au libelle ci-dessous.
\newcommand{\CAPETLIBELLEENONCES}{(\'enonc\'es)}

\ifdefined\CAPETMANUEL
  % --- Mode manuel : enonce en place, corrige differe vers l'annexe ---
  \newcommand{\CAPETSHOW}[2]{%
    \def\CAPETargun{#1}%
    \ifx\CAPETargun\CAPETLIBELLEENONCES
      #1%            appel decoratif : rendu tel quel, hors numerotation
    \else
      \CAPETSHOWEXO{#1}{#2}%
    \fi
  }
  \newcommand{\CAPETSHOWEXO}[2]{%
    \stepcounter{CAPETexo}%
    \csgdef{CAPETcorrige\theCAPETexo}{#2}%
    \edef\CAPETtmpnum{\theCAPETexo}%
    #1%
    \par\vspace{2pt}%
    {\footnotesize\itshape\color{colExo}%
     \hyperref[capet:cor:\CAPETtmpnum]{Corrige de l'exercice \CAPETtmpnum\ $\rightarrow$ p.~\pageref{capet:cor:\CAPETtmpnum}}%
     \label{capet:eno:\CAPETtmpnum}}%
  }
  % \CAPETANNEXE : rejoue tous les corriges captures, dans l'ordre, avec intertitres
  \newcommand{\CAPETANNEXE}{%
    \setcounter{CAPETannexe}{0}%
    \loop\ifnum\value{CAPETannexe}<\value{CAPETexo}%
      \stepcounter{CAPETannexe}%
      \edef\CAPETcurnum{\theCAPETannexe}%
      \ifcsname CAPETnotionstart\CAPETcurnum\endcsname
        \par\medskip
        {\normalsize\bfseries\color{colExo}\csuse{CAPETnotionstart\CAPETcurnum}}%
        \par\nobreak\smallskip
      \fi
      \begin{tcolorbox}[exoBox, breakable,
        title={Corrige de l'exercice \CAPETcurnum\
               \hfill\normalfont\footnotesize
               \hyperref[capet:eno:\CAPETcurnum]{enonce p.~\pageref{capet:eno:\CAPETcurnum}}}]%
      \label{capet:cor:\CAPETcurnum}%
      \csuse{CAPETcorrige\CAPETcurnum}%
      \end{tcolorbox}%
    \repeat
  }
\else
  % --- Modes historiques, comportement strictement inchange ---
  \newcommand{\CAPETSHOW}[2]{\ifdefined\CAPETCORRIGE #2\else #1\fi}
  \newcommand{\CAPETANNEXE}{}
\fi

% ── Difficulté 1 à 3 (étoiles pleines, amssymb) ───────────────────────
\newcommand{\difficulte}[1]{%
  \ifcase#1\relax
  \or $\bigstar$\;niveau 1/3%
  \or $\bigstar\bigstar$\;niveau 2/3%
  \or $\bigstar\bigstar\bigstar$\;niveau 3/3%
  \fi
}

% ── Style encadré extrait de document technique (banque d'exercices) ──
\tcbset{
  docTech/.style={
    breakable,
    enhanced,
    colback=colGrayLight,
    colframe=colSeparator,
    fonttitle=\bfseries\small\itshape,
    coltitle=black,
    colbacktitle=colGrayLight,
    left=4pt, right=4pt, top=3pt, bottom=3pt,
    before skip=4pt, after skip=4pt,
  }
}

% ── Style tcolorbox de section principale ────────────────────────────
% Refonte 2026 : bandeau titre coloré arrondi + fond très pâle, sans cadre.
\tcbset{
  ficheSection/.style={
    breakable, enhanced jigsaw,
    frame hidden, boxrule=0pt, arc=1.8mm,
    colback=#1!4!white,
    colbacktitle=#1,
    coltitle=white,
    fonttitle=\sffamily\bfseries,
    left=7pt, right=7pt, top=5pt, bottom=5pt,
    toptitle=3pt, bottomtitle=3pt,
    before skip=10pt, after skip=8pt,
  }
}

% ── Style encadré question type ──────────────────────────────────────
% Carte blanche, bandeau sombre, liseré très clair.
\tcbset{
  questiontype/.style={
    breakable, enhanced jigsaw,
    colback=white,
    frame hidden, boxrule=0pt, arc=1.8mm,
    colbacktitle=colElec!78!black,
    coltitle=white,
    fonttitle=\sffamily\bfseries\small,
    left=7pt, right=7pt, top=4pt, bottom=4pt,
    toptitle=2.5pt, bottomtitle=2.5pt,
    before skip=8pt, after skip=8pt,
    borderline={0.6pt}{0pt}{colElec!12!white},
  }
}

% ── Style encadré jury ───────────────────────────────────────────────
% Bandeau ambre, fond crème, sans cadre.
\tcbset{
  jury/.style={
    breakable, enhanced jigsaw,
    colback=colJuryBg,
    frame hidden, boxrule=0pt, arc=1.8mm,
    colbacktitle=colJuryFrame,
    coltitle=white,
    fonttitle=\sffamily\bfseries,
    left=7pt, right=7pt, top=4pt, bottom=4pt,
    toptitle=3pt, bottomtitle=3pt,
    before skip=10pt, after skip=8pt,
  }
}

% ── En-tête / pied de page communs (refonte 2026) ─────────────────────
% #1 code fiche, #2 titre court, #3 couleur domaine
\newcommand{\ficheHeaders}[3]{%
  \fancyhead[L]{\small\sffamily{\bfseries\textcolor{#3}{#1}}\hspace{0.7em}\textcolor{colSlate}{#2}}%
  \fancyhead[R]{\small\sffamily\textcolor{colSlate}{CAPET SII IE\, --- \,Session 2026}}%
  \fancyfoot[C]{\small\sffamily\textcolor{colSlate}{\thepage\,/\,\pageref{LastPage}}}%
  \fancyhead[C]{}\fancyfoot[L]{}\fancyfoot[R]{}%
}

% ── Bloc en-tête de fiche (BLOC 0, refonte 2026) ──────────────────────
% #1 code  #2 titre  #3 couleur domaine  #4 domaine  #5 ligne présence  #6 sous-thèmes
\newcommand{\ficheEntete}[6]{%
  \begin{tcolorbox}[enhanced, breakable,
    frame hidden, boxrule=0pt, arc=2.5mm,
    interior style={left color=#3, right color=#3!75!black},
    left=9pt, right=9pt, top=7pt, bottom=7pt,
    before skip=0pt, after skip=6pt]
    {\color{white}\sffamily\bfseries\LARGE #1\quad #2}\\[4pt]
    {\color{#3!22!white}\sffamily\small \MakeUppercase{#4}\quad\textbar\quad #5}\\[1pt]
    {\color{#3!35!white}\sffamily\footnotesize Sous-thèmes : #6}
  \end{tcolorbox}%
}

% ── Mise en page ─────────────────────────────────────────────────────
\pagestyle{fancy}
\fancyhf{}
\renewcommand{\headrulewidth}{0.5pt}
\renewcommand{\footrulewidth}{0pt}
\AtBeginDocument{%
  \renewcommand{\headrule}{{\color{colRule}\hrule width\headwidth height 0.5pt}}%
}

\setlist[itemize]{leftmargin=1.2em, itemsep=1pt, topsep=2pt, parsep=0pt}
\setlist[enumerate]{leftmargin=1.2em, itemsep=3pt, topsep=2pt, parsep=0pt}

\setlength{\parskip}{2pt}
\setlength{\columnsep}{10pt}

\def\CAPETBEGINFILE{\begin{document}}
\def\CAPETENDFILE{\end{document}}
\else
\def\CAPETBEGINFILE{}
\def\CAPETENDFILE{}
\fi

\ifdefined\var\else\DeclareSIUnit{\var}{var}\fi

\fancyhead[L]{\textbf{\textcolor{colExo}{T06}} - Exercices - Transformateur}
\fancyhead[R]{\small\textit{Banque d'exercices \CAPETSHOW{(\'enonc\'es)}{(corrig\'es)}}}
\fancyfoot[C]{\small\thepage\,/\,\pageref{LastPage}}

\CAPETBEGINFILE

\begin{tcolorbox}[
  enhanced, breakable,
  colback=colExo,
  colframe=colExo,
  left=6pt, right=6pt, top=5pt, bottom=5pt,
  before skip=0pt, after skip=4pt
]
  {\color{white}\bfseries\large T06 \quad Transformateur (monophas\'e et triphas\'e)}\\[2pt]
  {\color{white}\small Banque d'exercices -- 5 exercices \CAPETSHOW{(\'enonc\'es)}{(corrig\'es)}}
\end{tcolorbox}

\begin{tcolorbox}[exoBox,
  title={Exercice 1 -- \difficulte{1} \hfill\normalfont\footnotesize exercice in\'edit}
]
\CAPETSHOW{
Un transformateur triphas\'e de distribution $S_n = \SI{250}{\kilo\volt\ampere}$, $\SI{20000}{\volt}/\SI{400}{\volt}$ (couplage \'etoile-\'etoile), a pour caract\'eristiques d'essai en court-circuit : tension de court-circuit relative $\varepsilon_{cc} = \dfrac{U_{cc}}{U_{1n}} = 4\,\%$, pertes cuivre nominales (mesur\'ees \`a $I_{cc}=I_{2n}$) $P_{cc} = \SI{3,5}{\kilo\watt}$.

\begin{center}
\begin{circuitikz}[european, scale=0.85]
\draw
(0,0) node[left]{$V_{2n}$}
to[R, l=$R_s$] (2.5,0)
to[cute inductor, l=$X_s$] (5,0)
to[short, -*] (6,0) node[right]{$V_2$}
;
\draw (0,0) -- (0,-1.2) -- (6,-1.2) -- (6,0);
\node at (3,-1.5) {charge $\cos\varphi$};
\end{circuitikz}
\end{center}

\begin{enumerate}[leftmargin=1.5em]
  \item Calculer la composante r\'esistive relative $\varepsilon_R = P_{cc}/S_n$.
  \item En admettant l'approximation usuelle $\varepsilon_X \approx \varepsilon_{cc}$ (valable lorsque $R_s \ll X_s$, cas courant pour cette gamme de puissance), calculer la chute de tension relative $\Delta U/U_{2n} \approx \varepsilon_R\cos\varphi + \varepsilon_X\sin\varphi$ pour une charge \`a pleine puissance, $\cos\varphi = 0{,}8$ (inductif).
  \item En d\'eduire la tension secondaire en charge $U_2$, et commenter sa conformit\'e avec la tol\'erance usuelle de $\pm 5\,\%$ sur la tension de distribution basse tension.
\end{enumerate}

}{
\textbf{Corrig\'e.}

\begin{enumerate}[leftmargin=1.5em]
  \item $\varepsilon_R = \dfrac{P_{cc}}{S_n} = \dfrac{3\,500}{250\,000} = 0{,}014 = 1{,}4\,\%$.
  \item Avec $\varepsilon_X \approx \varepsilon_{cc} = 4\,\%$ et $\sin\varphi = \sqrt{1-0{,}8^2} = 0{,}6$ :
  \[ \dfrac{\Delta U}{U_{2n}} \approx \varepsilon_R\cos\varphi + \varepsilon_X\sin\varphi = 0{,}014\times0{,}8 + 0{,}04\times0{,}6 = 0{,}0112 + 0{,}024 = 0{,}0352 = 3{,}52\,\% \]
  \item $\Delta U = 0{,}0352 \times 400 \approx \SI{14,1}{\volt}$, donc $U_2 \approx 400 - 14{,}1 = \SI{385,9}{\volt}$.

  \textbf{Conclusion.} L'\'ecart relatif \`a $U_{2n}$ est de $3{,}52\,\%$, inf\'erieur \`a la tol\'erance usuelle de $\pm5\,\%$ admise sur la tension de distribution BT : le transformateur tient l'exigence de tension jusqu'\`a sa pleine puissance, sans qu'il soit n\'ecessaire d'agir sur le changeur de prises \`a vide.
\end{enumerate}
}
\end{tcolorbox}

% Exercice 2 -- docTech -- 2018, epreuve E2 (EDE ENE 2), questions 13 a 15.
% Remplace (passe complementarite 2026-09-02) l'ancien Ex.2 "P3 vs PD3 / rapport m"
% (2026 E1 partie E, questions E4 et E10) qui reproduisait a l'identique les QT1 et QT2 de la
% fiche T06 (memes valeurs 513 V, m = 2,23, 631 V, 446 V, m = 1,12) et empietait de surcroit
% sur la notion T08 (Redresseurs), qui traite la meme partie E questions E8 a E10.
% Tracabilite des donnees de l'encadre docTech (toutes textuellement presentes dans
% docs/.../06_PDFs_Originaux/2018/E2_Sujet_2018/E2_Sujet_2018.md) :
%   - Tableau "Tensions et Puissances pour le COMPACTO II(R)" : DRS 4 section 8, lignes 662 et
%     669 du .md (lignes kVA et A des gammes monophasee 230-240 V et triphasee 400-410 V).
%   - Lampe MASTER SON-T PIA Plus 150W : DRS 6, lignes 758-790 du .md (puissance nominale
%     154,0 W, courant lampe 1,8 A, tension 230 V, tension d'amorcage 2800 V, amorceur externe).
%   - Rendement > 98,5 % : DRS 4 section 9, ligne 700 du .md.
%   - Donnees d'enonce : 46 luminaires (l. 131), lampes SHP SON T PIA Plus 150 W avec amorceurs
%     Pa = 50 W et facteur de puissance 0,84 (l. 248-250), questions 13 a 15 (l. 255-259).
% Rien n'est invente : l'encadre ne fait que remettre en forme des valeurs textuellement
% presentes dans le .md du sujet. DRS 5 (releve sur 15 luminaires) est une image omise par la
% conversion Markdown : il n'est donc PAS reproduit, la question 1 s'appuie sur DRS 6 et DRS 4.
\begin{tcolorbox}[exoBox,
  title={Exercice 2 -- \difficulte{2} \hfill\normalfont\footnotesize 2018, EDE ENE 2, questions 13 \`a 15 (documents ressources DRS 4 et DRS 6)}
]
\CAPETSHOW{
Une municipalit\'e r\'enove l'\'eclairage public de la \guillemotleft\ rue des f\^etes \guillemotright, qui compte \textbf{46 luminaires} r\'epartis \'equitablement sur les trois phases d'un r\'eseau $230/\SI{400}{\volt}$. Chaque luminaire est constitu\'e d'une lampe au sodium haute pression SHP SON-T PIA Plus $\SI{150}{\watt}$ associ\'ee \`a un amorceur externe de puissance $P_a = \SI{50}{\watt}$ ; l'ensemble pr\'esente un facteur de puissance $\cos\varphi = 0{,}84$.

Pour abaisser la tension d'alimentation des lampes en nuit profonde, l'armoire est \'equip\'ee d'un \emph{r\'egulateur abaisseur de tension} COMPACTO II\textsuperscript{\textregistered} : un transformateur \`a rapport variable ins\'er\'e apr\`es le contacteur d'\'eclairage public. Il s'agit d'en choisir le calibre dans la gamme du constructeur.

\begin{center}
\begin{tikzpicture}[scale=1, >=Stealth, node distance=1.0cm]
  \node[draw, align=center, minimum width=2.0cm, minimum height=0.95cm] (res) {R\'eseau\\$230/\SI{400}{\volt}$};
  \node[draw, align=center, minimum width=1.7cm, minimum height=0.95cm, right=0.8cm of res] (kc) {Contacteur\\KC};
  \node[draw, align=center, minimum width=2.5cm, minimum height=0.95cm, right=0.8cm of kc] (reg) {R\'egulateur\\abaisseur $S_n = ?$};
  \node[draw, align=center, minimum width=2.3cm, minimum height=0.95cm, right=0.8cm of reg] (lum) {46 luminaires\\$\cos\varphi = 0{,}84$};
  \draw[-latex] (res) -- (kc);
  \draw[-latex] (kc) -- (reg);
  \draw[-latex] (reg) -- (lum);
\end{tikzpicture}
\end{center}

\begin{tcolorbox}[docTech, title={Document technique -- DRS 6 (lampe MASTER SON-T PIA Plus 150W) et DRS 4 \S 8 (gamme COMPACTO II\textsuperscript{\textregistered})}]
\small
\textbf{DRS 6 -- MASTER SON-T PIA Plus 150W/220 E40}, caract\'eristiques \'electriques :
\begin{center}
\begin{tabular}{ll}
\toprule
Tension d'alimentation de la lampe & $\SI{230}{\volt}$ \\
Puissance (valeur nominale) & $\SI{154,0}{\watt}$ \\
Courant lampe (EM) (nom.) & $\SI{1,8}{\ampere}$ \\
Tension d'amor\c{c}age (max.) & $\SI{2800}{\volt}$ \\
Description du syst\`eme & amorceur externe (E) \\
\bottomrule
\end{tabular}
\end{center}

\smallskip
\textbf{DRS 4 \S 8 -- Tensions et puissances pour le COMPACTO II\textsuperscript{\textregistered}} (rendement annonc\'e $> 98{,}5\,\%$) :
\begin{center}
\begin{tabular}{llrrrrr}
\toprule
\textbf{Monophas\'e 230--240 V} & (kVA) & 3 & 4 & 6 & 9 & 12 \\
                                & (A)   & 14 & 17 & 26 & 40 & 53 \\
\midrule
\textbf{Triphas\'e 400--410 V}  & (kVA) & 9 & 12 & 18 & 27 & 36 \\
                                & (A)   & 14 & 17 & 26 & 40 & 53 \\
\bottomrule
\end{tabular}
\end{center}
\end{tcolorbox}

\begin{enumerate}[leftmargin=1.5em]
  \item \textbf{Relever dans le document technique DRS 6} la puissance nominale r\'eellement absorb\'ee par la lampe seule, et la comparer aux $\SI{150}{\watt}$ de la d\'esignation commerciale. En d\'eduire la puissance active $P_{1L}$ d'un luminaire complet (lampe $+$ amorceur), puis calculer sa puissance r\'eactive $Q_{1L}$ et sa puissance apparente $S_{1L}$.
  \item Calculer la puissance apparente totale $S_{TL}$ appel\'ee par les 46 luminaires de la rue, puis la puissance apparente et le courant appel\'es \emph{par phase}.
  \item \`A l'aide du tableau DRS 4 \S 8, \textbf{choisir} le mod\`ele et le calibre du r\'egulateur abaisseur. Justifier le choix \`a la fois sur le crit\`ere de puissance apparente et sur le crit\`ere de courant, et pr\'eciser lequel des deux crit\`eres est ici le plus contraignant.
\end{enumerate}

}{
\textbf{Corrig\'e.}

\begin{enumerate}[leftmargin=1.5em]
  \item \emph{Lecture du DRS 6.} La d\'esignation commerciale annonce $\SI{150}{\watt}$, mais la ligne \guillemotleft\ Puissance (valeur nominale) \guillemotright\ du document constructeur indique $\SI{154,0}{\watt}$ : la lampe absorbe en r\'ealit\'e $2{,}7\,\%$ de plus que sa d\'esignation. Le dimensionnement du sujet retient la valeur de d\'esignation $P_{\text{lampe}} = \SI{150}{\watt}$ ; l'\'ecart de $\SI{4}{\watt}$ par lampe sera reconsid\'er\'e au moment de valider la marge du calibre choisi.
  \[ P_{1L} = P_{\text{lampe}} + P_a = 150 + 50 = \SI{200}{\watt} \]
  Avec $\cos\varphi = 0{,}84$, soit $\varphi = \arccos 0{,}84 \approx 32{,}9\,\text{\textdegree}$ et $\tan\varphi \approx 0{,}646$ :
  \[ S_{1L} = \frac{P_{1L}}{\cos\varphi} = \frac{200}{0{,}84} \approx \SI{238}{\volt\ampere}
     \qquad Q_{1L} = P_{1L}\tan\varphi \approx 200 \times 0{,}646 \approx \SI{129}{\var} \]
  \emph{V\'erification :} $\sqrt{200^2 + 129^2} \approx 238$, coh\'erent.
  \item Pour les 46 luminaires :
  \[ S_{TL} = 46 \times S_{1L} = 46 \times 238 \approx \SI{10,95}{\kilo\volt\ampere} \approx \SI{11,0}{\kilo\volt\ampere} \]
  R\'epartis \'equitablement sur trois phases, sous la tension simple $V = \SI{230}{\volt}$ :
  \[ S_{\text{phase}} = \frac{S_{TL}}{3} = \frac{10\,950}{3} \approx \SI{3,65}{\kilo\volt\ampere}
     \qquad I_{\text{phase}} = \frac{S_{\text{phase}}}{V} = \frac{3\,650}{230} \approx \SI{15,9}{\ampere} \]
  \item \emph{Lecture de la gamme DRS 4 \S 8.} Deux crit\`eres doivent \^etre satisfaits simultan\'ement.
  \begin{itemize}[leftmargin=1.2em]
    \item \emph{Crit\`ere en puissance apparente} : le r\'egulateur triphas\'e doit tenir $S_{TL} \approx \SI{11,0}{\kilo\volt\ampere}$. Dans la gamme triphas\'ee 400--410 V, le calibre $\SI{9}{\kilo\volt\ampere}$ est insuffisant ($9 < 11{,}0$) ; le calibre $\SI{12}{\kilo\volt\ampere}$ convient, avec une marge de $(12-11{,}0)/11{,}0 \approx 9\,\%$.
    \item \emph{Crit\`ere en courant} : chaque phase appelle $\SI{15,9}{\ampere}$. La ligne \guillemotleft\ (A) \guillemotright\ du tableau donne $\SI{14}{\ampere}$ pour le calibre $\SI{9}{\kilo\volt\ampere}$ -- insuffisant, le courant admissible serait d\'epass\'e de $14\,\%$ -- et $\SI{17}{\ampere}$ pour le calibre $\SI{12}{\kilo\volt\ampere}$, soit une marge de $(17-15{,}9)/15{,}9 \approx 7\,\%$.
  \end{itemize}
  On retient donc le \textbf{COMPACTO II\textsuperscript{\textregistered} triphas\'e 400--410 V de $\SI{12}{\kilo\volt\ampere}$, $\SI{17}{\ampere}$}, \'equivalent \`a trois modules monophas\'es de $\SI{4}{\kilo\volt\ampere}$ (m\^eme ligne de courant $\SI{17}{\ampere}$ dans la gamme monophas\'ee 230--240 V).

  Les deux crit\`eres \'ecartent le calibre inf\'erieur, mais c'est le \textbf{courant} qui est ici le plus contraignant : sa marge ($7\,\%$) est plus faible que celle en puissance apparente ($9\,\%$).
\end{enumerate}

\smallskip
\textbf{Conclusion.} La marge de $7\,\%$ en courant du calibre $\SI{12}{\kilo\volt\ampere}$ couvre pr\'ecis\'ement l'\'ecart relev\'e \`a la question 1 entre les $\SI{150}{\watt}$ commerciaux et les $\SI{154}{\watt}$ r\'eels du DRS 6 ($+2{,}7\,\%$) : le choix reste valide m\^eme en retenant la puissance constructeur, ce qui n'aurait pas \'et\'e le cas du calibre $\SI{9}{\kilo\volt\ampere}$. La le\c{c}on de dimensionnement est g\'en\'erale -- un transformateur, ou tout appareil \`a rapport variable qui en d\'erive, se choisit sur sa \emph{puissance apparente} et son \emph{courant admissible}, jamais sur la seule puissance active de la charge : ici l'\'ecart entre les deux atteint $19\,\%$ ($\SI{200}{\watt}$ contre $\SI{238}{\volt\ampere}$ par luminaire).
}
\end{tcolorbox}

\begin{tcolorbox}[exoBox,
  title={Exercice 3 -- \difficulte{2} \hfill\normalfont\footnotesize exercice in\'edit}
]
\CAPETSHOW{
Un transformateur triphas\'e peut \^etre coupl\'e \'etoile-\'etoile (Yy) ou \'etoile-triangle (Yd). Le couplage Yd introduit un d\'ephasage de $30^\circ$ entre tensions primaires et secondaires, et pr\'esente la propri\'et\'e de bloquer la circulation des courants harmoniques de rang multiple de 3 (rang 3, 9, 15...) entre les deux c\^ot\'es du transformateur.

\begin{center}
\begin{tikzpicture}[scale=0.9, >=Stealth]
  \node at (0,2.6) {\textbf{Primaire (\'etoile Y)}};
  \draw (0,2) -- (0,1) (0,1) -- (-0.87,0.5) (0,1) -- (0.87,0.5) (0,1) -- (0,0);
  \node[circ] at (0,1) {};
  \node[above] at (0,2) {$a$};
  \node[below left] at (-0.87,0.5) {$b$};
  \node[below right] at (0.87,0.5) {$c$};
  \node[below] at (0,0) {$N$};

  \node at (5,2.6) {\textbf{Secondaire (triangle d)}};
  \draw (4.2,0.3) -- (5.8,0.3) -- (5,1.8) -- (4.2,0.3);
  \node[below left] at (4.2,0.3) {$a'$};
  \node[below right] at (5.8,0.3) {$b'$};
  \node[above] at (5,1.8) {$c'$};
\end{tikzpicture}
\end{center}

\begin{enumerate}[leftmargin=1.5em]
  \item Rappeler la relation entre tension compos\'ee et tension simple pour un enroulement coupl\'e en \'etoile, puis pour un enroulement coupl\'e en triangle.
  \item Un redresseur en aval du transformateur g\'en\`ere des harmoniques de courant de rang 3 c\^ot\'e secondaire. Expliquer pourquoi un couplage Yd (avec triangle c\^ot\'e g\'en\'erateur d'harmoniques) emp\^eche ces harmoniques de rang 3 de remonter c\^ot\'e r\'eseau primaire.
  \item Quel est l'inconv\'enient pratique du couplage Yd par rapport au couplage Yy lorsqu'il s'agit de d\'eterminer le rapport de transformation $m$ \`a partir d'un cahier des charges exprim\'e en tensions secondaires ?
\end{enumerate}

}{
\textbf{Corrig\'e.}

\begin{enumerate}[leftmargin=1.5em]
  \item Couplage \'etoile : $U_{\text{compos\'ee}} = \sqrt{3}\,V_{\text{simple}}$ (les trois enroulements partagent un point neutre commun). Couplage triangle : chaque enroulement est directement soumis \`a la tension compos\'ee, donc $U_{\text{compos\'ee}} = V_{\text{enroulement}}$ (pas de facteur $\sqrt{3}$).
  \item Les courants harmoniques de rang 3 (et multiples de 3) sont en phase dans les trois enroulements (composante homopolaire). Dans un enroulement en triangle, ces courants homopolaires peuvent circuler en boucle ferm\'ee \`a l'int\'erieur du triangle sans emprunter les fils de ligne : ils restent donc confin\'es c\^ot\'e triangle et ne se retrouvent pas c\^ot\'e \'etoile (primaire), qui ne voit que les rangs non multiples de 3.
  \item En couplage Yd, le rapport $m$ relie une tension simple (c\^ot\'e \'etoile) \`a une tension compos\'ee (c\^ot\'e triangle) : il faut donc explicitement convertir les tensions du cahier des charges (souvent donn\'ees en tensions compos\'ees) avant d'appliquer $m$, sous peine d'une erreur de facteur $\sqrt{3}$. En Yy, cette conversion n'est pas n\'ecessaire car $m$ s'applique indiff\'eremment aux tensions simples ou compos\'ees (le facteur $\sqrt{3}$ se simplifie des deux c\^ot\'es).
\end{enumerate}

\textbf{Conclusion.} Le couplage n'est donc pas un simple d\'etail de c\^ablage : d\`es que la charge aval est polluante en harmoniques de rang 3 (ponts redresseurs, r\'ecepteurs monophas\'es r\'epartis sur les trois phases), le couplage Yd est impos\'e par l'exigence de qualit\'e de l'\'energie au point de livraison, au prix d'une conversion suppl\'ementaire en $\sqrt{3}$ dans le calcul de $m$ et d'un indice horaire ($30^\circ$) \`a respecter pour toute mise en parall\`ele avec un autre transformateur.
}
\end{tcolorbox}

% Exercice 4 -- inedit -- dimensionnement du circuit magnetique (formule de Boucherot).
% Remplace (passe complementarite 2026-09-02) l'ancien Ex.4 "essai a vide, calcul de Rf et Xm"
% (2026 E2 BAC+3, questions A12-A13) qui reproduisait a l'identique la QT3 de la fiche T06
% (meme source, memes donnees V = 230 V, I0 = 42,5 A, P0 = 7,82 kW, memes resultats
% Rf = 20,3 ohms et Xm = 5,62 ohms). Le corpus de sujets 2013-2026 ne contient que deux etudes
% de transformateur (2026 E1 partie E, deja QT1/QT2 et reservee a T08 ; 2024 E1 partie B, deja
% utilisee par T06 ex. 2 anterieurement et surtout par T03 ex. 2) : le remplacant est donc un
% exercice INEDIT, conformement au dernier recours prevu par le protocole.
% Il porte sur un travail intellectuel absent des trois QT de la fiche : le dimensionnement du
% circuit magnetique (nombre de spires a partir de l'induction admissible), et non le calcul
% d'un rapport de transformation ni l'identification d'un modele a vide.
% Ordres de grandeur verifies (WebSearch 2026-09-02) : formule de Boucherot V = 4,44 N f Bmax S
% (le coefficient 4,44 vaut pi*racine(2)) ; les aciers au silicium a grains orientes utilises
% dans les transformateurs de distribution saturent vers 2,0 T, d'ou une induction
% d'exploitation usuelle de 1,6 a 1,7 T (fr.wikipedia.org/wiki/Formule_de_Boucherot ;
% donnees de saturation des toles a grains orientes).
% Le transformateur 250 kVA 20 kV/400 V est celui des exercices 1 et 5 de cette meme notion,
% pour assurer la coherence interne du fichier.
\begin{tcolorbox}[exoBox,
  title={Exercice 4 -- \difficulte{2} \hfill\normalfont\footnotesize exercice in\'edit}
]
\CAPETSHOW{
On reprend le transformateur triphas\'e de distribution des exercices 1 et 5 : $S_n = \SI{250}{\kilo\volt\ampere}$, $\SI{20000}{\volt}/\SI{400}{\volt}$, couplage \'etoile-\'etoile, $f = \SI{50}{\hertz}$. Il s'agit maintenant d'en dimensionner le \emph{circuit magn\'etique}, c'est-\`a-dire de d\'eterminer le nombre de spires de chaque enroulement.

La loi de Faraday appliqu\'ee \`a un enroulement de $N$ spires embrassant un flux sinuso\"idal d'amplitude $\Phi_{max} = B_{max}\,S$ dans un noyau de section $S$ conduit \`a la \textbf{formule de Boucherot} :
\[ V = \frac{2\pi}{\sqrt{2}}\,N\,f\,B_{max}\,S = 4{,}44\,N\,f\,B_{max}\,S \]
o\`u $V$ est la valeur efficace de la tension aux bornes de l'enroulement.

Le constructeur emploie des t\^oles d'acier au silicium \`a grains orient\'es, dont l'induction de saturation avoisine $\SI{2,0}{\tesla}$. Pour \'eviter tout risque de saturation en r\'egime nominal, il impose une induction maximale d'exploitation $B_{max} \le \SI{1,6}{\tesla}$. La section utile d'un noyau vaut $S = \SI{150}{\centi\metre\squared}$.

\begin{center}
\begin{tikzpicture}[scale=0.85, >=Stealth]
  % noyau trois colonnes
  \draw[line width=2.2pt, colSlate] (0,0) rectangle (7,4);
  \draw[line width=2.2pt, colSlate] (3.5,0) -- (3.5,4);
  % enroulements schematises sur la colonne centrale
  \foreach \y in {1.1,1.5,1.9,2.3,2.7} {
    \draw[colElec, line width=1pt] (3.15,\y) arc (270:90:0.18);
    \draw[colExo, line width=1pt] (3.85,\y) arc (270:450:0.18);
  }
  \node[colElec, left]  at (2.85,2.6) {\footnotesize $N_1$ (HTA)};
  \node[colExo,  right] at (4.15,2.6) {\footnotesize $N_2$ (BT)};
  \draw[<->, thick] (0.55,4.35) -- (2.95,4.35);
  \node[above] at (1.75,4.3) {\footnotesize colonne, section $S$};
  \draw[->, thick] (5.6,2.6) -- (5.6,1.4);
  \node[right] at (5.65,2.0) {\footnotesize $\Phi(t)$, $B_{max}$};
\end{tikzpicture}
\end{center}

\begin{enumerate}[leftmargin=1.5em]
  \item Le couplage \'etant \'etoile des deux c\^ot\'es, exprimer la tension $V_2$ aux bornes d'un enroulement secondaire (donc d'une colonne) en fonction de la tension compos\'ee secondaire $U_2 = \SI{400}{\volt}$, et la calculer. Faire de m\^eme pour $V_1$ au primaire.
  \item En d\'eduire le nombre de spires secondaire $N_2$ strictement n\'ecessaire pour ne pas d\'epasser $B_{max} = \SI{1,6}{\tesla}$. Le nombre de spires \'etant entier, indiquer -- en le justifiant physiquement -- s'il faut arrondir \`a l'entier inf\'erieur ou sup\'erieur, puis donner la valeur retenue et l'induction maximale r\'eellement atteinte.
  \item En d\'eduire le nombre de spires primaire $N_1$, et v\'erifier que l'induction calcul\'ee c\^ot\'e primaire est bien la m\^eme (les deux enroulements embrassent le m\^eme flux).
  \item Calculer le courant nominal secondaire par phase. Sachant que la densit\'e de courant admissible dans les conducteurs de cuivre refroidis par bain d'huile est de l'ordre de $J = \SI{3}{\ampere\per\milli\metre\squared}$, en d\'eduire la section de cuivre d'un conducteur secondaire.
\end{enumerate}

}{
\textbf{Corrig\'e.}

\begin{enumerate}[leftmargin=1.5em]
  \item En couplage \'etoile, chaque enroulement (une colonne) est soumis \`a la \emph{tension simple} :
  \[ V_2 = \frac{U_2}{\sqrt{3}} = \frac{400}{\sqrt{3}} \approx \SI{231}{\volt}
     \qquad V_1 = \frac{U_1}{\sqrt{3}} = \frac{20\,000}{\sqrt{3}} \approx \SI{11547}{\volt} \]
  C'est la premi\`ere source d'erreur du calcul : appliquer Boucherot avec la tension compos\'ee $\SI{400}{\volt}$ conduirait \`a un nombre de spires surestim\'e d'un facteur $\sqrt{3}$.
  \item Avec $S = \SI{150}{\centi\metre\squared} = \SI{150e-4}{\metre\squared}$ :
  \[ N_2 = \frac{V_2}{4{,}44\,f\,B_{max}\,S} = \frac{230{,}9}{4{,}44 \times 50 \times 1{,}6 \times 150\times10^{-4}} = \frac{230{,}9}{5{,}328} \approx 43{,}3 \text{ spires} \]
  \emph{Sens de l'arrondi.} \`A tension et section impos\'ees, $B_{max}$ varie en $1/N_2$ : \emph{diminuer} le nombre de spires \emph{augmente} l'induction. Arrondir \`a 43 spires ferait donc d\'epasser la limite de $\SI{1,6}{\tesla}$. Il faut arrondir \`a l'entier \textbf{sup\'erieur} :
  \[ \boxed{N_2 = 44 \text{ spires}} \qquad
     B_{max,\text{r\'eel}} = \frac{V_2}{4{,}44\,f\,N_2\,S} = \frac{230{,}9}{4{,}44 \times 50 \times 44 \times 150\times10^{-4}} \approx \SI{1,58}{\tesla} \]
  soit une marge de $\SI{0,42}{\tesla}$ (environ $21\,\%$) vis-\`a-vis de la saturation \`a $\SI{2,0}{\tesla}$.
  \item Le rapport de transformation, identique en tensions simples ou compos\'ees puisque les deux c\^ot\'es sont en \'etoile :
  \[ m = \frac{V_2}{V_1} = \frac{U_2}{U_1} = \frac{400}{20\,000} = 0{,}02
     \qquad N_1 = \frac{N_2}{m} = \frac{44}{0{,}02} = \boxed{2\,200 \text{ spires}} \]
  \emph{V\'erification c\^ot\'e primaire :}
  \[ B_{max} = \frac{V_1}{4{,}44\,f\,N_1\,S} = \frac{11\,547}{4{,}44 \times 50 \times 2\,200 \times 150\times10^{-4}} \approx \SI{1,58}{\tesla} \]
  On retrouve bien la m\^eme valeur : les deux enroulements sont travers\'es par le m\^eme flux, l'induction est une propri\'et\'e du \emph{noyau}, pas de l'enroulement.
  \item Courant nominal secondaire (ligne, donc \'egal au courant d'enroulement en \'etoile) :
  \[ I_{2n} = \frac{S_n}{\sqrt{3}\,U_2} = \frac{250\,000}{\sqrt{3} \times 400} \approx \SI{361}{\ampere} \]
  \[ s_{cu} = \frac{I_{2n}}{J} = \frac{361}{3} \approx \SI{120}{\milli\metre\squared} \]
  soit, pour un conducteur m\'eplat, une section de l'ordre de $\SI{120}{\milli\metre\squared}$ par spire secondaire -- d'o\`u le recours \`a des barres de cuivre m\'eplates et non \`a du fil rond \`a ce niveau de courant.
\end{enumerate}

\smallskip
\textbf{Conclusion.} Le dimensionnement retenu ($N_1 = 2\,200$, $N_2 = 44$, $S = \SI{150}{\centi\metre\squared}$) conduit \`a $B_{max} \approx \SI{1,58}{\tesla}$, en de\c{c}\`a de la limite d'exploitation de $\SI{1,6}{\tesla}$ fix\'ee par le constructeur : le crit\`ere de non-saturation est satisfait. Cette marge n'est pas un luxe -- la norme admet une surtension permanente de $10\,\%$ sur le r\'eseau de distribution, qui porterait l'induction \`a $1{,}58 \times 1{,}1 \approx \SI{1,74}{\tesla}$, encore inf\'erieure aux $\SI{2,0}{\tesla}$ de saturation. Arrondir $N_2$ \`a 43 spires aurait au contraire donn\'e $B_{max} \approx \SI{1,61}{\tesla}$ en r\'egime nominal et $\SI{1,77}{\tesla}$ \`a $+10\,\%$, avec un courant magn\'etisant fortement d\'eform\'e par l'approche de la saturation.
}
\end{tcolorbox}

\begin{tcolorbox}[exoBox,
  title={Exercice 5 -- \difficulte{3} \hfill\normalfont\footnotesize exercice in\'edit}
]
\CAPETSHOW{
Le transformateur de l'exercice 1 ($S_n = \SI{250}{\kilo\volt\ampere}$, $P_{cc} = \SI{3,5}{\kilo\watt}$ \`a pleine charge) a des pertes fer constantes $P_{fer} = \SI{800}{\watt}$ (ind\'ependantes de la charge, li\'ees uniquement \`a $V_1$). On note $k$ le taux de charge ($k=1$ \`a pleine charge), de sorte que les pertes cuivre valent $P_{cu}(k) = k^2 P_{cc}$.

\begin{center}
\begin{tikzpicture}[scale=0.9, >=Stealth]
  \draw[-latex, thick] (0,0) -- (2.5,0) node[midway, above]{$P_1$};
  \node[draw, minimum height=1.2cm, minimum width=1cm] at (3,0) {};
  \draw[-latex, thick] (3.5,0) -- (6,0) node[midway, above]{$P_2$};
  \draw[-latex, thick] (3,-0.6) -- (3,-1.8) node[midway, right]{$P_{fer}+P_{cu}(k)$};
\end{tikzpicture}
\end{center}

\begin{enumerate}[leftmargin=1.5em]
  \item Exprimer le rendement $\eta(k)$ en fonction de $k$, $S_n$, $\cos\varphi$, $P_{fer}$ et $P_{cc}$.
  \item Calculer $\eta$ \`a pleine charge ($k=1$), pour $\cos\varphi = 0{,}8$.
  \item Montrer que le rendement est maximal lorsque $P_{fer} = P_{cu}(k)$, en d\'eduire le taux de charge optimal $k_{opt}$ et le rendement maximal correspondant.
\end{enumerate}

}{
\textbf{Corrig\'e.}

\begin{enumerate}[leftmargin=1.5em]
  \item Puissance utile $P_2 = k\,S_n\cos\varphi$, pertes totales $P_{fer} + k^2 P_{cc}$ :
  \[ \eta(k) = \dfrac{k\,S_n\cos\varphi}{k\,S_n\cos\varphi + P_{fer} + k^2 P_{cc}} \]
  \item \`A $k=1$ : $P_2 = 250\,000\times0{,}8 = \SI{200}{\kilo\watt}$, pertes $= 800+3\,500 = \SI{4,3}{\kilo\watt}$.
  \[ \eta(1) = \dfrac{200\,000}{200\,000+4\,300} = \dfrac{200\,000}{204\,300} \approx 0{,}9789 = 97{,}89\,\% \]
  \item $\eta(k)$ est maximal lorsque $\dfrac{\mathrm{d}\eta}{\mathrm{d}k}=0$, ce qui \'equivaut (apr\`es calcul, num\'erateur de la d\'eriv\'ee nul) \`a $P_{fer} = k^2 P_{cc}$, soit pertes fixes = pertes variables. D'o\`u :
  \[ k_{opt} = \sqrt{\dfrac{P_{fer}}{P_{cc}}} = \sqrt{\dfrac{800}{3\,500}} = \sqrt{0{,}2286} \approx 0{,}478 \]
  Soit environ $47{,}8\,\%$ de la charge nominale. \`A ce point : $P_2 = 0{,}478\times250\,000\times0{,}8 \approx \SI{95,6}{\kilo\watt}$, pertes totales $= 2\times800 = \SI{1,6}{\kilo\watt}$ :
  \[ \eta_{max} = \dfrac{95\,600}{95\,600+1\,600} = \dfrac{95\,600}{97\,200} \approx 0{,}9835 = 98{,}35\,\% \]
\end{enumerate}

\textbf{Conclusion.} Le rendement maximal est obtenu \`a $k_{opt} \approx 48\,\%$ de charge, et non \`a pleine charge : c'est ce qui justifie le \emph{surdimensionnement} volontaire des transformateurs de distribution, choisis pour que leur taux de charge moyen r\'eel (de l'ordre de $50\,\%$) co\"incide avec l'optimum de rendement -- crit\`ere repris par les exigences d'\'eco-conception europ\'eennes sur les pertes des transformateurs de puissance.
}
\end{tcolorbox}

\CAPETENDFILE
